% ============================================================
% Appendix D -- pointer only.
%
% The working itself moved to the compendium (chapter 19) on 3 September 2026:
% ten pages of alternative aggregations was more length than a robustness check
% should cost the thesis. This stub stays because four places cite
% \ref{app:aggregation}, one of them a generated file, and because a reader who
% wants the check should be told in one sentence where it is and what it says.
% ============================================================

\chapter{How the Reported Ratios Were Aggregated}\label{app:aggregation-chapter}
\label{app:aggregation}

A ratio is a nonlinear function of two means, so it depends on the order in
which the averaging is done: $\E[X]/\E[Y]$ is not $\E[X/Y]$, and the gap
between them grows with the spread. Both ratios reported in this thesis
therefore need the same question asked of them, namely whether the number is a
property of the estimator or of the estimand chosen to summarise it.

Both were recomputed under four alternative aggregations of the same frozen
cells, refitting nothing. The outcome, in one sentence each. For the
unstructured comparison the direction is the same under every aggregation and
the estimand reported in Chapter~\ref{ch:em-results} is systematically the
largest of them, which is stated wherever that number appears. For the
structure-aware comparison the reported estimand is \emph{not} uniformly the
largest, but the floor across all four aggregations is
$\structaggfloor\times$ with a spread of $\structaggspread$, so the reading of
a low single-digit multiple at the smallest training set growing to roughly six
at the largest survives every one of them.

The full recomputation, both tables and the protocol behind them, is
Chapter~19 of \compendium{}. It is generated by
\codefile{tools/make\_aggregation\_robustness.py} from the same frozen outputs
that produce the headline table, so the two cannot drift apart.
